A unital ring $B$ is an \emph{exchange ring} if, for every
$b\in B$, there is an idempotent $p\in Bb$ with
$1-p\in B(1-b)$. This is the ordinary exchange condition
\cite[Proposition~1.1 and Theorem~2.1]{NicholsonExchange}; it
includes semiperfect rings. It permits a stronger object-lifting
statement than the Jacobson case alone.

\begin{theorem}[coefficient quotients with exchange upper rings]
\label{thm:exchange-coefficient-lifting}
Let $T=(Q,F)$ be any unital inverse semiring such that every
$F(e)$, $e\in U(Q)$, is an exchange ring. Let $\mathcal I$
be any compatible coefficient ideal, and let $T/\mathcal I$
retain the support $Q$.

For every finite idempotent support matrix $E$, every idempotent
of $A_E(T/\mathcal I)$ lifts to an idempotent of $A_E(T)$.
Thus the lift has exactly the prescribed support, even when
that additive matrix fiber is nonunital. Every finitely generated
projective over $T/\mathcal I$ lifts, and the maps
\[
 V(T)\longrightarrow V(T/\mathcal I),\qquad
 K_0(T)\longrightarrow K_0(T/\mathcal I)
\]
are surjective. For the compatible Jacobson ideal $C(T)$ of
Theorem~\ref{thm:upper-jacobson-quotient}, both maps are
isomorphisms.
\end{theorem}

No nilpotence or radical condition is imposed on $\mathcal I$.
Injectivity for an arbitrary coefficient ideal is not asserted.
The local-idempotent multiplicity classification in
Theorem~\ref{thm:upper-semiperfect-amplitudes} retains its
semiperfectness hypothesis.
